\icode{E-MP-SI-L2-260627-001}
氢原子光谱在可见光区域内有四条谱线\(\mathrm{H}_{\alpha}\)、\(\mathrm{H}_{\beta}\)、\(\mathrm{H}_{\gamma}\)和\(\mathrm{H}_{\delta}\)，都是氢原子中电子从量子数\(n > 2\)的能级跃迁到\(n = 2\)的能级时发出的光，它们在真空中的波长由长到短，可以判定\paren[]
\begin{choices}
    \item \(\mathrm{H}_{\alpha}\)对应的前后能级之差最小
    \item 同一介质对\(\mathrm{H}_{\alpha}\)的折射率最大
    \item 同一介质中\(\mathrm{H}_{\delta}\)的传播速度最大
    \item 用\(\mathrm{H}_{\gamma}\)照射某一金属能发生光电效应，则\(\mathrm{H}_{\beta}\)也一定能
\end{choices}
\begin{answer}
    A\qquad 提示：\(\mathrm{H}_{\alpha}\)频率\(\nu\)最小，折射率最小\(n = \frac{c}{v}\)
\end{answer}